Operator learning constructs fast surrogates for the solution maps of
parametric partial differential equations and physical forward
models.\footnote{Code, data pointers, and the per-run summaries behind every
number reported here are at
\url{https://github.com/yspennstate/neural-means-kernel-corrections}.}
Neural operators \citep{kovachki2023neural,lu2021deeponet,li2021fourier}
learn their own representations; kernel and Gaussian-process methods
\citep{batlle2024kernel,mora2025operator,owhadi2019operator} come with exact
solves and error certificates and a strong record on the same benchmarks.
This paper combines the two: a neural network mean, trained in the metric the
application reports, with an exact Mat\'ern kernel regression of its residual
and of its learned features. We evaluate the combination in depth on two
public problems with published baselines --- the structural-mechanics
benchmark of \citet{dehoop2022cost} and the OCO-2 radiative-transfer
emulation problem of \citet{susiluoto2025radiative} --- and, with one
unchanged recipe, across the seven-benchmark operator-learning suite
assembled by \citet{batlle2024kernel} from the problems of
\citet{dehoop2022cost} and \citet{lu2022comprehensive}.
Table~\ref{tab:suite} collects the published numbers and ours.

\begin{table}[t]
\centering
\caption{The two problems studied in depth (OCO-2, structural mechanics) and
the rest of the operator-learning suite, with the best published mean
relative $L^2$ test error and ours. OCO-2 rows compare against the
Gaussian-process emulator of \citet{susiluoto2025radiative} on its own test
points, in that paper's reduced-coefficient metric; Section~\ref{sec:oco2}
reports both of its metrics, on which one surrogate wins both on each of the
three bands. The
published suite numbers are per-problem methods tuned by their authors; ours
outside the two problems studied in depth come from one recipe run unchanged
(a Mat\'ern kernel at the median length scale, fit capped at 6000 points, a
residual MLP mean, outputs reduced by PCA, the stage selected on validation).
Burgers and Darcy use slightly smaller training splits than their baselines
(800 and 824 pairs) and Darcy a different grid, from the copies we could
obtain.}
\label{tab:suite}
\begin{tabular}{llcc}
\toprule
Problem & Map & Best published & This work \\
\midrule
OCO-2, O2 band & atmospheric state $\to$ radiance & 16.89\% & $4.12\pm0.05\%$ \\
OCO-2, weak CO$_2$ & atmospheric state $\to$ radiance & 24.06\% & $16.28\pm0.06\%$ \\
OCO-2, strong CO$_2$ & atmospheric state $\to$ radiance & 16.14\% & $8.08\pm0.01\%$ \\
Structural mechanics & boundary load $\to$ stress & 4.55\% (PARA-Net) & $\errPipe\pm\errPipeSd\%$ \\
\midrule
Burgers & initial condition $\to$ solution & 1.93\% (FNO) & 2.38\% \\
Darcy flow & permeability $\to$ pressure & 2.32\% (POD-DeepONet) & 2.01\% \\
Advection~I & smooth pulse $\to$ solution & $\approx0\%$ (kernel) & --- \\
Advection~II & discontinuous pulse $\to$ solution & 11.28\% (linear kernel) & 11.29\% \\
Helmholtz & wavespeed $\to$ field & 1.00\% (kernel) & 1.75\% \\
Navier--Stokes & vorticity $\to$ vorticity & 0.12\% (kernel) & 0.51\% \\
\bottomrule
\end{tabular}
\end{table}

The two problems studied in depth sit at opposite ends of the
neural-versus-kernel spectrum, and the pair is the point. On structural
mechanics the families tie: the pipeline reaches $\errPipe\pm\errPipeSd\%$
under the 20000-sample protocol, replicated at ten independently seeded runs,
which is below every published method but PARA-Net's \pubBestHi\% and within one
standard error of that one; and \errBestLow\% under the 1250-sample protocol
against a published best of \pubBestLo\%. What the coupling adds there is mostly
understanding, and the replication is what licenses the reading. The residuals
of every architecture we train correlate between \corrNetLo{} and \corrNetHi{}
at every seed; their shared component concentrates where the finite element
data are least reliable and is flat in ensemble diversity and training-set
size; the hardest cases are hard for all members at once, so averaging does not
touch them; and the spread of our own pipeline across seeds is $\errPipeSd$
points, some sixty times smaller than the $0.65$ points that separate the
published architectures. The plateau near 4.5\% therefore reads as a property
of the benchmark's data rather than of any surrogate class, and it is not
run-to-run noise that makes it look flat (Section~\ref{sec:diversity}). Absent
regenerated finite-element data this remains an inference, but it is the reading
every measurement we made supports. On the rest of the suite the same recipe sorts the problems by the
same logic: where the map is smooth and data are plentiful (Helmholtz,
Navier--Stokes, Darcy) the kernel stage wins the pipeline's internal
selection and lands within a factor of two of the tuned per-problem kernels
of \citet{batlle2024kernel}, and on Burgers the corrected mean sits near the
Fourier neural operator.

On OCO-2 the same components separate by an order of magnitude, and the
pipeline's job changes: not to couple two comparable models but to move the
kernel's exactness onto the network's representation. Per spectral band of
the satellite, the task maps a reduced atmospheric state to a reduced
radiance spectrum, and the baseline is the kernel-flow emulator of
\citet{susiluoto2025radiative}, whose stored predictions we score on the same
test points. An exact Mat\'ern head on the trained network's features reaches
$4.1\%$ where the same kernel on the raw state reaches $40\%$; the effective
dimension of the two Gram matrices is unchanged, the target's squared
native-space norm falls by a factor of about forty, and
Proposition~\ref{prop:pullback} gives the mechanism. Trained in a loss
weighted towards the reported metric and combined per output coordinate, the
result improves on the emulator on both of that problem's error metrics at
once on all three bands, at ten of ten seeds on two of them and nine of ten
on the third (Section~\ref{sec:oco2} is precise about the margins).

Each stage of the pipeline carries one supporting result, and the objects the
results cover are the deployed estimators rather than idealizations of them:
a symmetrization lemma behind the reflection averaging; a second-moment
identity that predicts stacking outcomes from measured residual correlations,
sharpened to a two-sided bracket that certifies, from the same measurements,
how close any convex ensemble must sit to its floor; finite-sample oracle
inequalities for the validated simplex stack, for the per-pixel affine stack
the pipeline actually fits, and for the per-coordinate selection that builds
the OCO-2 surrogate, chained through the grid-tuned correction into a
certificate for the shipped object whose every constant is measured and
released; an optimal-recovery bound $\|G-m\|_K\,P_\lambda(u)$ for the
corrected surrogate; a pullback identity and a rate for the feature kernel;
and a distribution-free coverage statement whose hypotheses the experiment
satisfies exactly, the calibration set being disjoint from every selection.
The coverage statement matters because the design factor $P_\lambda$, though
a valid bound, turns out to be a poor pointwise ranking of the actual error,
and its split-conformal rescaling is the only uncertainty signal that
survives our tests (Section~\ref{sec:uq}). Every stated result is also
checked numerically: identities to machine precision, inequalities on
constructions whose ground truth is known exactly, and the pipeline-level
statements on the released artifacts of every seed
(Appendix~\ref{app:verification}). The two-member stacking rule in
particular is scored prospectively, on every member pair of every seed rather
than on selected examples, and that scoring is reported as it came out: the
rule is reliable outside a margin band set by its own estimation noise and no
better than chance inside it, which is a limitation of the measurement we can
state precisely rather than one we can remove
(Section~\ref{sec:secmom-margin}). The symmetrization inequality, which is
a statement about an expectation, is reported with the one finite-sample
reversal that ten seeds turn up.

Section~\ref{sec:benchmark} describes the problems, protocols, and published
results. Section~\ref{sec:method} specifies the pipeline.
Sections~\ref{sec:experiments} and~\ref{sec:oco2} report the two studies,
including the data-scaling experiments. Section~\ref{sec:theory} states the
supporting theory, one result per stage, with proofs in
Appendix~\ref{app:proofs} and their numerical verification in
Appendix~\ref{app:verification}, and Section~\ref{sec:discussion} discusses
limitations.
